\documentclass[11pt,a4paper]{article}
\usepackage[margin=27mm]{geometry}
\usepackage[T1]{fontenc}
\usepackage{lmodern,microtype,amsmath,amssymb,amsthm,booktabs,array}
\usepackage[hidelinks]{hyperref}
\hypersetup{pdftitle={An Ant on the Clock: Exponential Escape Bounds for Langton's Ant},pdfauthor={Hao Ke}}
\usepackage{fancyhdr}
\pagestyle{fancy}\fancyhf{}\fancyhead[L]{\small An Ant on the Clock}\fancyfoot[C]{\thepage}
\setlength{\headheight}{14pt}
\newtheorem{theorem}{Theorem}[section]
\newtheorem{lemma}[theorem]{Lemma}
\newtheorem{corollary}[theorem]{Corollary}
\theoremstyle{remark}\newtheorem{remark}[theorem]{Remark}
\newcommand{\Z}{\mathbb Z}\newcommand{\N}{\mathbb N}\newcommand{\Q}{\mathbb Q}
\newcommand{\R}{\mathbb R}\newcommand{\PhiC}{\Psi_{a,b,c}}
\newcommand{\eps}{\varepsilon}\newcommand{\symdiff}{\mathbin{\triangle}}
\title{An Ant on the Clock:\\Exponential Escape Bounds for Langton's Ant}
\author{Hao Ke\\Independent researcher\\\texttt{i@kehao.me}}
\date{15 October 2026}
\begin{document}
\maketitle
\begin{abstract}
We prove that the classical two-colour Langton ant exits every square of side
$n$ within $(13/4)3^n+4$ updates, uniformly over all initial colourings and
internal poses. For a rectangle with $N$ columns and fixed height $H$, we
obtain an upper bound of order $N^{\lfloor H/2\rfloor+1}$. A family of
rational potentials on the critical curve $(a-1)(b-1)=4$ has positive monomial
increments. Four-step costs give the square estimate, while pairing along the
short direction gives the rectangular bound. A related strict rational clock
yields decidable complete-state reachability, two-sided orbit equivalence and
a computable orbit normal form. Exact polynomial checks, finite-state audits
and Lean proofs of the supercritical local clock lemmas accompany the arguments.
\end{abstract}

\section{Model and results}
A state is $X=(B,p,h)$, where $B\subset\Z^2$ is finite, $p\in\Z^2$, and
$h\in\{N,E,S,W\}$. Headings are encoded cyclically by $0,1,2,3$.
Let $e_h$ be the corresponding unit vector. A state at time $t$ is observed
\emph{before} update $t$. The literal update is
\begin{equation}\label{eq:rule}
 h'=h+1-2\mathbf1_B(p)\pmod4,\qquad
 B'=B\symdiff\{p\},\qquad p'=p+e_{h'}.
\end{equation}
Thus white gives a right turn, black a left turn, and the read cell flips
before the ant moves.

Let $S(N,H)$ be the maximum first-exit time from
$Q_{N,H}=\{0,\ldots,N-1\}\times\{0,\ldots,H-1\}$ over all internal
colourings, positions and headings. The exit move is counted. Exterior colours
can be arbitrary: no exterior cell is read before the first exit.

\begin{theorem}\label{crit:thm:main}
For positive integers $n,N,H$, put
\[
 R_n=3^{1-n}(3^n-1)^2+4(3^{n-1}-3^{1-n}).
\]
Then
\begin{align}
 S(n,n)&\le4\left\lfloor\frac{3R_n}{16}\right\rfloor+4
       \le\frac{13}{4}3^n+4,\label{crit:eq:square}\\
 S(N,H)&\le\left\lfloor\frac92(NH+16)(4N+1)^{\lfloor H/2\rfloor}\right\rfloor+3.
 \label{crit:eq:rectangle}
\end{align}
In particular, $S(N,H)=O_H(N^{\lfloor H/2\rfloor+1})$.
\end{theorem}

The square estimate improves the factorial upper bound of
Etse~\cite{Etse2026} to exponential order in the side length. His fixed-height
rectangle bound has degree $H-1$ for even $H$ and $H-2$ for odd $H$; the degree
above is smaller for every $H\ge6$. For heights two and three, his linear
estimates remain sharper.

The mechanism is a rational potential with positive local increments. Its
range on a finite domain bounds the time spent there. At critical parameters,
individual increments can be arbitrarily small, but short blocks have a uniform
positive cost. Isotropic and anisotropic weights then give the two bounds.
For a specified complete target state, a supercritical member of the clock
family gives a finite search budget and a canonical representative of each
orbit. These complete-state questions complement the circuit-prediction and
infinite-background constructions of~\cite{Gajardo2002}.

We first give the strict clock and its elementary increment identity, then
combine critical components to prove the square and rectangle bounds. Finally,
we use the fixed strict clock for complete-state orbit algorithms.

\section{An explicit strict rational clock}
Define
\[
 \chi(x,y)=(-1)^{x+y},\qquad
 \sigma(X)=\begin{cases}\chi(p),&h=N,S,\\-\chi(p),&h=E,W,\end{cases}
\]
and put $\eps(h)=+1$ for $h=N,E$, and $\eps(h)=-1$ for $h=S,W$.
Each update changes both the lattice parity and the heading axis, so $\sigma$
is invariant under $F$.

The construction uses the diagonal coordinate $x+y$, which changes by
exactly $+1$ or $-1$ on every move. We assign symmetric exponential weights
to the two directions along this coordinate. A colour flip changes one term
of the weighted field sum; a heading-dependent pose term compensates for that
change so that their combined increment is positive in every turn case.
The invariant $\sigma$ aligns the field signs with the ant's heading axis.

Fix $c\in\Z$ and a rational $r>3$. The parameter $c$ centres the diagonal
coordinate, while $r$ controls the growth of the weights. Write
$s(q)=q_x+q_y-c$ and set
\begin{align}
 f_r(s)&=r^s+r^{-s},\\
 g_r(s,+1)&=\frac{2}{r+1}(r^s-r^{1-s}),\\
 g_r(s,-1)&=\frac{2}{r+1}(r^{-s}-r^{s+1}),\\
 \Phi_{r,c}(X)&=-\sigma(X)\sum_{q\in B}\chi(q)f_r(s(q))
                    +g_r(s(p),\eps(h)).\label{eq:clock}
\end{align}
Finite support makes \eqref{eq:clock} a computable rational number, including
when coordinates are negative.

\begin{lemma}[Increment identity]\label{lem:increment}
Let $U=r^{s(p)}$, $V=r^{-s(p)}$, and
\[
 d_r=\frac{r-3}{r+1},\qquad E_r=\frac{5r+1}{r+1}.
\]
If $e=\eps(h)$ and $e'=\eps(h')$, then
\[
\begin{array}{c|c}
 (e,e')&\Phi_{r,c}(FX)-\Phi_{r,c}(X)\\\hline
 (+,+)\text{ or }(-,-)&d_r(U+V)\\
 (+,-)&d_rU+E_rV\\
 (-,+)&E_rU+d_rV
\end{array}
\]
In particular $\Phi_{r,c}(FX)-\Phi_{r,c}(X)\ge2d_r>0$.
\end{lemma}
\begin{proof}
Let $\delta=+1$ on a white read and $\delta=-1$ on a black read. From the
four headings and two colours in \eqref{eq:rule},
\[
 \delta\sigma(X)\chi(p)=ee',\qquad s(p')=s(p)+e'.
\]
The black-field contribution to the increment is therefore
$-ee'(U+V)$. Subtracting the old pose term from the new one gives
\[
 -ee'(U+V)+g_r(s(p)+e',e')-g_r(s(p),e).
\]
Substitution of the two definitions of $g_r$ gives the displayed table.
For example, for $(e,e')=(+,+)$, the pose increment is
$2(r-1)(U+V)/(r+1)$, which yields $d_r(U+V)$ after the field increment.
For completeness, after multiplying by $r+1$, the other three cases are
\begin{align*}
(-,-):\quad&-(r+1)(U+V)+2(rV-U)-2(V-rU),\\
(+,-):\quad& (r+1)(U+V)+2(rV-U)-2(U-rV),\\
(-,+):\quad& (r+1)(U+V)+2(rU-V)-2(V-rU).
\end{align*}
Expansion gives the corresponding polynomial identities. Since $UV=1$, $U+V\ge2$;
also $E_r\ge d_r>0$. These give the lower bound in every case.
\end{proof}

For complete-state orbit algorithms, one fixed member
suffices: choosing $r=4$ and $c=0$ gives the following uniform time estimate.

\begin{theorem}[Strict clock]\label{thm:clock}
The function $\Phi=\Phi_{4,0}$ is computable on every complete finite-support
state and satisfies
\begin{equation}\label{eq:time}
 \Phi(F^tX)-\Phi(X)\ge\frac25t\qquad(t\in\N).
\end{equation}
\end{theorem}
\begin{proof}
Apply Lemma~\ref{lem:increment} with $r=4$ and sum the increments.
\end{proof}

\section{The critical four-quadrant potential}
Use the same checkerboard invariant $\sigma$ as above. Fix $a,b>1$ and a centre $c=(c_x,c_y)\in\R^2$, and put
$U(q)=a^{q_x-c_x}$, $V(q)=b^{q_y-c_y}$. Define
\begin{equation}\label{crit:eq:potential}
 \PhiC(X)=-\sigma(X)\sum_{q\in B}\chi(q)
 (U(q)-U(q)^{-1})(V(q)-V(q)^{-1})+G_h(U(p),V(p)),
\end{equation}
where
\begin{align*}
 G_N(U,V)&=\frac2{b+1}(U+U^{-1})(V-bV^{-1}),\\
 G_E(U,V)&=\frac2{a+1}(U-aU^{-1})(V+V^{-1}),\\
 G_S(U,V)&=\frac2{b+1}(U+U^{-1})(V^{-1}-bV),\\
 G_W(U,V)&=\frac2{a+1}(U^{-1}-aU)(V+V^{-1}).
\end{align*}
This finite sum keeps the entire black field, including old debris. It is
rational for rational $a,b$ and integer $c$.

\begin{lemma}[Critical identity]\label{crit:lem:identity}
Suppose $(a-1)(b-1)=4$. At a state with $U=U(p)$ and $V=V(p)$, the increment
$\PhiC(FX)-\PhiC(X)$ is given by Table~\ref{crit:tab:increments}. In particular it
is strictly positive for every complete finite-support state.
\end{lemma}
\begin{table}[ht]
\centering
\begin{tabular}{ccc}\toprule
Heading before update & White/right & Black/left\\\midrule
$N$ & $4U/V$ & $4/(UV)$\\
$E$ & $4/(UV)$ & $4V/U$\\
$S$ & $4V/U$ & $4UV$\\
$W$ & $4UV$ & $4U/V$\\\bottomrule
\end{tabular}
\caption{Exact increments on the critical curve.}\label{crit:tab:increments}
\end{table}
\begin{proof}
Let $\delta=1$ on white and $-1$ on black, and let $\eta_h=1$ on vertical
headings and $-1$ on horizontal headings. Only the current colour changes, so
the field contribution is
$-\eta_h\delta(U-U^{-1})(V-V^{-1})$.
If $(d_x,d_y)$ is the outgoing unit vector and $h'$ the new heading, add
$G_{h'}(Ua^{d_x},Vb^{d_y})-G_h(U,V)$. In each of the eight cases, its difference
from the stated table entry simplifies to
\[
 \frac{(U^2+1)(V^2+1)(ab-a-b-3)}{UV(a+1)(b+1)}.
\]
The numerator vanishes on the critical curve.
\end{proof}

The potential combines rotated critical components of the directional
exponential family. A single component can have zero increment. Their sum does
not, but its increments need not have a location-independent lower bound.
The next sections derive such lower bounds for short blocks instead.

\section{A uniform four-step cost and square escape}
\begin{lemma}\label{crit:lem:four}
Let $a=b=3$ and let $c_x,c_y$ both be integers or both be half-integers. Then
\[
 \PhiC(F^4X)-\PhiC(X)\ge16/3
\]
for every complete finite-support state $X$.
\end{lemma}
\begin{proof}
Each entry in Table~\ref{crit:tab:increments}, divided by four, is an integer power of~3.
A nonconstant four-letter turn word contains consecutive opposite turns. For
such a pair the reward sum is
$4(3^j+3^{-j-1})$, with $j\in\Z$. For example, an $N$-facing right turn followed
by a left turn contributes $4U/V+4V/(3U)$. The other pairs are read directly
from the table. For integer $j$, $3^j+3^{-j-1}\ge1+1/3$: the expression increases
for $j\ge0$ and is symmetric under $j\mapsto-j-1$. The remaining increments
are positive.

If all four turns agree, they trace a unit square. For $N$-facing $RRRR$, the
sum is
\[
 4\left(U+\frac1{3U}\right)\left(\frac1V+\frac V3\right)\ge16/3
\]
by two applications of AM--GM. For each of the eight constant-turn/heading cases, the product of the four
positive rewards is $4^4/81$: the $U,V$ exponents cancel. Four-term AM--GM
therefore gives the same bound for all eight cases.
\end{proof}

The bound is attained: centre at $(1/2,1/2)$, colour $\{0,1\}^2$ black, and start
at $(0,0)$ facing south. Four left turns erase the square and increase the
potential by exactly $16/3$. This is not an optimality assertion for escape.

We next bound the range on $Q_{n,n}$. Set $c_x=c_y=(n-1)/2$, $s=x+y-(n-1)$,
$t=x-y$, and
\[
 f(k)=3^k+3^{-k},\qquad
 g(k,+)=\frac{3^k-3^{1-k}}2,\quad
 g(k,-)=\frac{3^{-k}-3^{k+1}}2.
\]
With $\epsilon(h)=+$ for $h=N,E$ and $-$ for $h=S,W$, elementary expansion of
\eqref{crit:eq:potential} gives
\begin{equation}\label{crit:eq:diagonal}
 \PhiC(X)=\sigma\sum_{q\in B}\chi(q)(f(t(q))-f(s(q)))
       +g(s(p),\epsilon(h))+g(t(p),\epsilon(h-1)).
\end{equation}
Here $h-1$ uses the indexing $N,E,S,W=0,1,2,3$. Since $s,t$ are integers,
\eqref{crit:eq:diagonal} is rational even when $n$ is even.

\begin{lemma}\label{crit:lem:range}
With fixed exterior colouring and fixed $\sigma$, the range of $\PhiC$ over
states whose heads lie in $Q_{n,n}$ is at most $R_n$ from Theorem~\ref{crit:thm:main}.
\end{lemma}
\begin{proof}
For a finite complete field, the exterior contribution is constant. For arbitrary
exterior colouring, sum only over the square; the same update identity holds
between internal states because every changed cell lies inside. Thus the
argument uses a finite sum in either case. For each diagonal,
\[
 \sum_{q\in Q_{n,n}} f(s(q))
   =2\,3^{1-n}\left(\frac{3^n-1}{2}\right)^2,
\]
and the same identity holds for $t$. The range of a binary weighted sum is
at most the sum of coefficient magnitudes. Using
$|f(t)-f(s)|\le f(t)+f(s)$ bounds the field range by
$3^{1-n}(3^n-1)^2$.
For $|k|\le n-1$, the combined range of $g(k,+)$ and $g(k,-)$ is at most
$2(3^{n-1}-3^{1-n})$, by monotonicity of the displayed powers.
The two head terms have range at most twice this, proving the claim.
\end{proof}

\begin{proof}[Proof of \eqref{crit:eq:square}]
If $T$ is the first exit, there are $T-1$ internal-to-internal updates.
Group them into disjoint four-step blocks. Lemmas~\ref{crit:lem:four}
and~\ref{crit:lem:range}, and positivity of leftover increments, imply
\[
 \left\lfloor\frac{T-1}{4}\right\rfloor\le
 \left\lfloor\frac{3R_n}{16}\right\rfloor.
\]
An infinite confined orbit is also excluded by applying this estimate to
arbitrarily long internal prefixes. Finally,
$R_n=13\,3^{n-1}-6-3^{2-n}<13\,3^{n-1}$ yields the coarser displayed bound.
\end{proof}

\begin{corollary}
For any fixed integer centre, the rational function
$\Theta(X)=\sum_{i=0}^3\Psi_{3,3,c}(F^iX)$ satisfies
$\Theta(FX)-\Theta(X)\ge16/3$.
\end{corollary}
\begin{proof}
The difference telescopes to the four-step increment. Computing the three
successors requires no infinite-future information.
\end{proof}
Either this finite-lookahead clock or the supercritical clock can support
complete-state orbit algorithms. We use the simpler fixed clock of
Theorem~\ref{thm:clock} in Section~\ref{sec:orbit}.

\section{Pairing in the short direction}
For $Q_{N,H}$ choose
\[
 a=1+1/N,\qquad b=1+4N,\qquad
 c=(0,\lfloor(H-1)/2\rfloor),\qquad k=\lfloor H/2\rfloor.
\]
Then $(a-1)(b-1)=4$, $|y-c_y|\le k$, and $1\le U\le a^{N-1}<3$.
For the last inequality, the binomial expansion bounds
$(1+1/N)^N$ by $\sum_{j=0}^N1/j!<3$.

\begin{lemma}\label{crit:lem:rectangle}
While the head remains in $Q_{N,H}$, a valid range bound is
\[
 R_{N,H}=3(NH+16)b^k.
\]
A two-update segment beginning with a horizontal heading, whose two read
positions are internal, increases the potential by at least
\[
 \delta_N=4a^{-(N-1)}(1+1/b)>4/3.
\]
\end{lemma}
\begin{proof}
Each field coefficient has magnitude at most $3b^k$, since
$|U-U^{-1}|<3$ and $|V-V^{-1}|\le b^k$. The field range is therefore at most
$3NHb^k$. Each $G_h$ is a sum of four monomials with coefficient magnitudes at
most~2; each monomial has magnitude at most $3b^k$. Hence
$|G_h|\le24b^k$, and the head range is at most $48b^k$.

The first move in the asserted pair is vertical, say by $v\in\{-1,1\}$.
The two entries of Table~\ref{crit:tab:increments} have vertical factors
$b^j$ and $b^{-j-1}$, where $j=v(y-c_y)$ is an integer.
Both are read at the same $x$, so both horizontal factors are at least
$a^{-(N-1)}$. The elementary integer inequality
$b^j+b^{-j-1}\ge1+1/b$ proves the asserted cost. It is greater than $4/3$
because $a^{N-1}<3$.
\end{proof}

\begin{proof}[Proof of \eqref{crit:eq:rectangle}]
The arrival axes alternate. The internal-to-internal updates can thus be grouped
into the pairs of Lemma~\ref{crit:lem:rectangle}, omitting at most one initial and
one final update. With $r$ pairs, $T\le2r+3$ and $r\delta_N\le R_{N,H}$.
An indefinitely confined trajectory contradicts the same finite range.
Consequently
\[
 T\le\left\lfloor\frac{2R_{N,H}}{\delta_N}\right\rfloor+3
 \le\left\lfloor\frac92(NH+16)(4N+1)^k\right\rfloor+3.
\]
\end{proof}

The proof also gives a sharper computable bound by replacing $R_{N,H}$ with
\[
 \widehat R_{N,H}=\sum_{q\in Q_{N,H}}|(U(q)-U(q)^{-1})(V(q)-V(q)^{-1})|
 +\max_{q,h}G_h(q)-\min_{q,h}G_h(q).
\]
The bound $\lfloor2\widehat R_{N,H}/\delta_N\rfloor+3$ uses only finite rational
arithmetic. It is not the exact maximum first-exit time.

By a quarter-turn rotation one can exchange $N$ and $H$ and take the
smaller of the two bounds. Etse's Theorem~12~\cite{Etse2026}, for $N\ge H\ge2$,
has polynomial degrees $H-1$ and $H-2$ in the even and odd cases respectively.
The degree $\lfloor H/2\rfloor+1$ is strictly smaller for $H\ge6$.

\section{Complete-state orbit algorithms}\label{sec:orbit}
Fix the single clock $\Phi=\Phi_{4,0}$ throughout this section.
Here the comparison is between complete states, so the same fixed clock is
used for every input. Its uniform increment converts a target clock value
into a finite simulation budget.

\begin{theorem}[Forward reachability]\label{thm:reach}
Given complete finite-support states $X,Y$, it is decidable whether
$F^tX=Y$ for some $t\ge0$. If such a time exists it is unique and satisfies
\[
 t\le N(X,Y):=\left\lfloor\frac52(\Phi(Y)-\Phi(X))\right\rfloor.
\]
\end{theorem}
\begin{proof}
Different $\sigma$ values or a negative clock difference rule out a match.
Otherwise compute the displayed integer and simulate through times $0$ to
$N(X,Y)$, inclusive, comparing the entire black set, position and heading.
Inequality~\eqref{eq:time} rules out every later match. Strict increase also
rules out two different matching times. Each state at finite time is finite,
and all rational arithmetic and full-state comparisons are effective.
\end{proof}

The inverse update is equally explicit. Given $(B,p,h)$, put $q=p-e_h$.
If $q\in B$, the preceding update was white-to-black, so the old heading was
$h-1$. If $q\notin B$, it was black-to-white, so the old heading was $h+1$.
In either case flip $q$ and put the predecessor at $q$. This defines a computable
bijection on the set $\mathcal X$ of complete finite-support states.

\begin{corollary}[Orbit equivalence]
The relation $\exists t\in\Z\;F^tX=Y$ is decidable, with a computable signed
time offset whenever it holds.
\end{corollary}
\begin{proof}
If the states differ, their time order on an orbit must agree with the strict
order of their clock values. Apply Theorem~\ref{thm:reach} in that order;
equal clock values allow a match only for identical states.
\end{proof}

\begin{theorem}[Computable orbit section]\label{thm:section}
Define the decidable set
\[
 \mathcal S=\{C\in\mathcal X:\Phi(C)\ge0>\Phi(F^{-1}C)\}.
\]
Every two-sided orbit meets $\mathcal S$ exactly once. Its representative $C(X)$
and the unique integer $k(X)$ with $X=F^{k(X)}C(X)$ are computable. The map
$X\mapsto(C(X),k(X))$ is a computable bijection between $\mathcal X$ and
$\mathcal S\times\Z$ under which $F$ acts by $(C,k)\mapsto(C,k+1)$.
\end{theorem}
\begin{proof}
Forward values tend to $+\infty$ by \eqref{eq:time}. Applying the same inequality
to $F^{-t}X$ shows that backward values tend to $-\infty$. Strict increase
therefore yields exactly one crossing from negative to nonnegative values.
Search forward when $\Phi(X)<0$ and backward otherwise, counting the signed
number of moves. A coarse bound for the moves and final predecessor check is
$\lceil(5/2)|\Phi(X)|\rceil+1$. The inverse of the displayed correspondence
computes $F^kC$ using forward or inverse updates as appropriate.
\end{proof}

The normal form separates the position along an orbit from its orbit label
$C\in\mathcal S$. A row-parity invariant shows that infinitely many distinct
labels are needed.

\begin{lemma}[Infinitely many orbit representatives]
The section $\mathcal S$ is infinite.
\end{lemma}
\begin{proof}
For each integer row $j$, the bit
\[
 I_j(X)=|B\cap(\Z\times\{j\})|
        +\mathbf1_{\{p_y=j,\ h\in\{E,W\}\}}\pmod2
\]
is invariant. On a vertical-to-horizontal turn the new head stays in the
read row and its indicator cancels that row's colour flip. On a
horizontal-to-vertical turn the old head indicator cancels the flip.
All other rows are unchanged. States with the head at the origin facing north
and sole black cell $(0,j)$ have distinct invariant vectors $(I_i)_{i\in\Z}$.
They therefore lie on distinct orbits, each with its own section representative.
\end{proof}

\paragraph{Computational cost.}
The clock gives explicit finite simulation budgets for the orbit algorithms.
With sparse binary coordinates, exact arithmetic may require exponential
space in the coordinate bit length: the expanded rational representation of
$4^x$ has length proportional to $|x|$. The algorithms establish decidability
with this representation-dependent cost.

\section{Verification and reproducibility}
The escape bounds combine exact local identities, short-block costs and finite
geometric-series range estimates. The supplementary artifact supplies two independently written
steppers and exact rational arithmetic. The retained baseline replay checks
4,704 local cases, 18,696 states in squares through side three, and 100 bound
parameters. A new integer-vector implementation checks 9,720 local cases,
including negative coordinates, remote black cells, and bases arbitrarily close
to three. A separate bit-mask functional graph exhausts all 4,213,000
colour/pose states of squares of sides one through four. Their maximum exit
times are $1,6,17,46$. Every graph transition is evaluated once, cycles are
explicitly rejected, and a maximizing witness for each size is replayed with
the integer-vector stepper. Direct field sums are compared with the geometric
formula in 60 cases; the coarse bound is checked at 1,000 side lengths.
These retained checks support the supercritical clock and small-square dynamics.

The artifact also includes kernel-checked Lean 4.30.0 proofs of all four
polynomial increment identities after clearing the positive denominator,
$U+V\ge2$ when $U>0$ and $UV=1$, the scaled increment lower bound, and the
finite telescoping/interval argument. Their printed dependencies are only
\texttt{propext}, \texttt{Classical.choice}, and \texttt{Quot.sound}.
The formalized components are the local algebra and telescoping lemmas.
The connection to the literal ant and the range estimates are analytic. The
critical eight-case identity and the short-block arguments are checked by exact
coefficient arithmetic and the written proofs; the Lean lemmas cover the
supercritical construction.

For the critical potential, the retained exact checker exercises 3,000 local
identity cases, 6,400 four-step cases, 10,608 horizontal-arrival pairs and
144 rectangular ranges. Its separate bit simulation exhausts 28,808 initial
states over eight small rectangles; 80 additional trajectories retain exterior
debris. A further implementation imports none of that code and checks the eight
cleared-denominator identities as Laurent polynomials in $a,b,U,V$. It also
verifies all eight opposite-turn pairs, all eight constant four-turn products,
the eight rectangular-pair cases and exit-step rounding. These algebraic checks
cover the local parameters symbolically; the universal range and coverage
arguments are the proofs in Sections~3--5.

\section{Conclusion}
Rational potentials connect local ant updates to quantitative control of
trajectories. At critical parameters, short-block costs remove the extra
factor of $n$ from the earlier square estimate and yield polynomial escape
bounds whose fixed-height degree is $\lfloor H/2\rfloor+1$. The same clock
framework also gives finite search budgets for complete-state targets and a
unique zero-crossing representative for each two-sided orbit. The geometric
estimates and the orbit algorithms use two forms of the same principle:
positive clock growth converts a finite range into a finite time budget.

\paragraph{Source and software availability.}
The sources, retained certificates and verification reports are in the
\href{https://github.com/kehao95/langtons-ant-core/tree/main/clock_escape}{public clock project}.
From the repository root, run:
\begin{verbatim}
python3 clock_escape/reproduce.py
python3 clock_escape/reproduce.py --critical-only
make -C clock_escape/paper
\end{verbatim}
The first command reproduces the complete evidence; the second runs only the
critical strengthening with Python's standard library. Source manifests record
SHA256 digests and the exact scope of the formalized components.

\paragraph{AI assistance.}
Generative AI was used in the research conversation to propose and critique
arguments, prepare proof and code drafts, and organize this manuscript.

\begin{thebibliography}{9}
\bibitem{Etse2026}
Kossi Roland Etse.
\newblock How Long Can The Escaping Ant Be Confined?
\newblock In \emph{51st International Symposium on Mathematical Foundations of
Computer Science (MFCS 2026)}, LIPIcs 386, 49:1--49:17, 2026.
\newblock \href{https://doi.org/10.4230/LIPIcs.MFCS.2026.49}{doi:10.4230/LIPIcs.MFCS.2026.49}.

\bibitem{Gajardo2002}
Anah\'{\i} Gajardo, Andr\'{e}s Moreira and Eric Goles.
\newblock Complexity of Langton's ant.
\newblock \emph{Discrete Applied Mathematics} 117:41--50, 2002.
\newblock \href{https://doi.org/10.1016/S0166-218X(00)00334-6}{doi:10.1016/S0166-218X(00)00334-6}.
\end{thebibliography}
\end{document}
